% Independent refraction arithmetic; unit vectors, outward-facing normal.
\clearpage
\section*{Verification case: distinguish total internal reflection from rounding}
\noindent In the existing optics derivation, the incident unit direction
\(\mathbf i\) points toward an interface and the surface unit normal
\(\mathbf n\) points against incidence. Define
\(c_i=-\mathbf i\cdot\mathbf n\) and \(\eta=n_1/n_2\).
The transmitted direction, if it exists, is
\[
\mathbf t=\eta\mathbf i+(\eta c_i-\sqrt{k})\mathbf n,
\qquad k=1-\eta^2(1-c_i^2).
\]
Suppose a ray passes from refractive index \(1.5\) to \(1.0\)
at \(45^\circ\), with \(\mathbf i\cdot\mathbf n=-1/\sqrt2\).
Then
\[
k=1-1.5^2\left(1-\frac12\right)=-0.125.
\]
There is no real transmitted ray: this is \emph{total internal
reflection}. Returning a unit transmitted direction by forcing
\(k\) to zero would produce a geometrically plausible but
physically incorrect result.

At \(30^\circ\) incidence from the same medium, however,
\(k=1-2.25(1-\tfrac34)=0.4375\).
The transmitted sine is \(0.75\), so the refraction angle is
approximately \(48.59^\circ\). The computed transmitted vector
should have unit length (within rounding) and point into the
second medium. The critical angle is
\(\arcsin(1/1.5)\), approximately \(41.81^\circ\).

\paragraph{Independent checks.}
Test both directions of propagation, critical incidence, slightly
subcritical and supercritical incidence, and input normal reversal
under a \emph{declared} convention. Reject zero-length normals and
nonfinite input. For \(k<0\), distinguish a tiny roundoff interval
from physically meaningful total internal reflection using a
well-justified numerical policy. Do not confuse this directional
calculation with Fresnel splitting, absorption or spectral dispersion.
